\chapter{Scaling Analysis: Impact of Cohort Size on Predictability}
\label{ch:scale}

\section{Scaling Analysis Objectives}

Stratifying 400 training subscribers across four behavioral user groups leaves smaller clusters (such as $G2$, comprising 14\% of the cohort) with approximately 56 training trajectories. Estimating robust subgroup transition probabilities from limited samples presents empirical bounds.

This observation leads to the central research question evaluated in this chapter: \emph{how does increasing subscriber cohort size impact next-cell prediction accuracy, and at what scale does performance saturate?}

\paragraph{Data Provenance.} Evaluation spans three dataset cohorts from \textbf{Table~\ref{tab:datasets}}: 400 subscribers (\Dmob{}, 100 held-out subscribers, baseline 43.1\%), 900 to 1,900 subscribers (\Dtwok{}, 100 held-out subscribers, baseline 48.2\%), and 5,000 to 21,000 subscribers (\Deleven{}, \textbf{1,000 held-out subscribers}, baseline 50.2\%). Because baseline scores vary across cohorts, evaluations report accuracy deltas relative to baseline benchmarks rather than absolute percentages.

\section{Experimental Setup for Cohort Size Sweep}

The cohort scaling evaluation sweeps training population size across eight nested steps (400 to 21,000 subscribers) while holding model architecture, LoRA adapter configuration, hour tokens, and per-group routing constant. Nested sampling ensures that larger cohorts strictly expand smaller cohorts without introducing population composition shifts. Across training sizes, training log counts increase from 369,451 to 1,538,379 events.

% F14 -- The whole ladder: gap to each run's own baseline, against training-set size.
\begin{figure}[t]
\centering
\plotlink{\nbfinal}{\fitwidth{%
\begin{tikzpicture}
\begin{axis}[rep, width=0.97\linewidth, height=6.4cm,
  xmode=log, log basis x=10,
  xmin=380, xmax=30000,
  xtick={400,900,1400,1900,5000,10000,15000,21000},
  xticklabels={400,900,1{,}400,1{,}900,5{,}000,10{,}000,15{,}000,21{,}000},
  x tick label style={font=\scriptsize,rotate=40,anchor=east},
  ymin=-0.5, ymax=7.6, ytick={0,1,...,5},
  yticklabels={0,$+1$,$+2$,$+3$,$+4$,$+5$},
  xlabel={training users (logarithmic scale)},
  ylabel={points gained over that run's own baseline},
  xmajorgrids, ymajorgrids,
]
% regime bands
\fill[clnonsig!30] (axis cs:380,-0.5) rectangle (axis cs:452,7.6);
\fill[clgain!12]   (axis cs:452,-0.5) rectangle (axis cs:1900,7.6);
\fill[clbase!12]   (axis cs:1900,-0.5) rectangle (axis cs:30000,7.6);
\node[font=\scriptsize,text=clrule,anchor=north,align=center] at (axis cs:900,7.45)
  {\textbf{GENERALISATION}\\ the four profiles become\\ large enough to read\\ \textcolor{clgain}{$+3.6\pt$}};
\node[font=\scriptsize,text=clrule,anchor=north,align=center] at (axis cs:8000,7.45)
  {\textbf{SATURATION}\\ every new user is a\\ behavioural duplicate\\ \textcolor{clbase}{$+0.1\pt$}};
\node[font=\scriptsize,text=clrule,anchor=south,align=center,rotate=90] at (axis cs:414,1.2)
  {\textbf{IMITATION}};
%
\addplot[clmodel,thick,mark=*,mark size=1.8pt] coordinates {
(400,0.6)(900,2.8)(1400,3.5)(1900,4.2)(5000,4.3)(10000,4.4)(15000,4.7)(21000,4.3)};
\addplot[clbase,dashed,thick,domain=380:30000,samples=2]{0};
%
\node[font=\scriptsize,anchor=west] at (axis cs:432,0.95) {$+0.6$};
\node[font=\scriptsize,anchor=north] at (axis cs:900,2.62) {$+2.8$};
\node[font=\scriptsize,anchor=north] at (axis cs:1400,3.32) {$+3.5$};
\node[font=\scriptsize,anchor=north west] at (axis cs:1950,4.15) {$+4.2$};
\node[font=\scriptsize,anchor=south] at (axis cs:5000,4.35) {$+4.3$};
\node[font=\scriptsize,anchor=south] at (axis cs:10000,4.45) {$+4.4$};
\node[font=\scriptsize,anchor=south] at (axis cs:15000,4.75) {$+4.7$};
\node[font=\scriptsize,anchor=south] at (axis cs:21000,4.35) {$+4.3$};
\node[font=\scriptsize,text=clbase,anchor=east] at (axis cs:28000,0.24) {the persistence baseline of each run};
\end{axis}
\end{tikzpicture}}}
\caption{The whole scaling ladder, in points over each run's own baseline}
\label{fig:scaling-ladder}
\figtext{%
Eight training-set sizes, each read against the persistence baseline measured
on its \emph{own} test users. The horizontal axis is logarithmic, so equal
distances are equal multiples and 400 to 1{,}900 occupies the same width as
4{,}000 to 19{,}000; the vertical axis is not accuracy but the gap to the
baseline, the only quantity that survives a change of test set, so the dashed
line at zero means ``indistinguishable from repeating the last cell''. At 400
training users the model gains $+0.6\pt$ and is barely distinguishable from the rule it
is meant to beat. Between 400 and 1{,}900 the edge climbs $+0.6 \rightarrow
+2.8 \rightarrow +3.5 \rightarrow +4.2\pt$, a four-fold increase in population
buying $+3.6\pt$, which is essentially the entire benefit the project ever
obtained from data volume. The ten-fold increase that follows buys $+0.1\pt$,
and not monotonically: the best tier (15{,}000) is followed by the worst
(21{,}000), and none of the steps inside that stretch is statistically
distinguishable from zero (Table~\ref{tab:bootstrap}). The curve does not slow
down, it stops. Two things constrain how far that reads. The 900 to
1{,}900-user tiers are scored on \stwok{} (100 test users, 48.2\,\% baseline)
and the 400-user point is the \smob{} run (43.1\,\% baseline) and the 5{,}000 to
21{,}000-user tiers on \seleven{} (1{,}000 test users, 50.2\,\% baseline),
which is why the axis is a gap and why the small step across that boundary is
not a measurement. And every tier received the same budget of optimiser steps,
so at 21{,}000 users only 36\,\% of the available training windows were ever
read: the defensible claim is that 1{,}900 users are already enough at this
training budget, not that more users can never help.}
\figsource{\nbfinal}{train\_cellid\_final.ipynb}{the eight-tier scaling notebook}
\end{figure}


\textbf{Figure~\ref{fig:scaling-ladder}} plots prediction accuracy delta relative to baseline across the eight training cohort sizes. The empirical curve exhibits two distinct regimes: rapid performance gains up to approximately 2,000 subscribers (+4.2 percentage points over baseline in \textbf{Figure~\ref{fig:scaling-ladder}}), followed by performance saturation across larger training cohorts up to 21,000 subscribers (+4.3 percentage points in \textbf{Figure~\ref{fig:scaling-ladder}}).

\subsection{Statistical Validation of Performance Saturation}

To establish whether performance plateauing above 2,000 subscribers reflects true saturation rather than random variance, paired bootstrap resampling \textbf{\cite{efron1993bootstrap}} was conducted across the 1,000 held-out subscribers.

% T9 -- Paired bootstrap over the test users (slide 75).
\begin{table}[!htbp]
\centering\small
\caption{Is any step up the ladder real? Every interval says no}
\label{tab:bootstrap}
\begin{tabular}{@{}lrlc@{}}
\toprule
\textbf{Step in training people} & \textbf{Measured difference} & \textbf{Range it could really be} & \textbf{Real gain?} \\
\midrule
5{,}000 $\rightarrow$ 10{,}000  & $+0.20\pt$ & $-0.21 \ \ldots \ +0.64$ & no \\
10{,}000 $\rightarrow$ 15{,}000 & $+0.09\pt$ & $-0.32 \ \ldots \ +0.52$ & no \\
15{,}000 $\rightarrow$ 21{,}000 & $-0.25\pt$ & $-0.72 \ \ldots \ +0.20$ & no \\
\midrule
5{,}000 $\rightarrow$ 21{,}000 \textit{(total)} & $+0.04\pt$ & $-0.43 \ \ldots \ +0.51$ & no \\
\bottomrule
\end{tabular}

\figtext{%
\figpara{The question.}{Between two training sizes the score moves a little. Is
that movement the model genuinely improving, or is it the luck of which thousand
people happened to be in the test set?}

\figpara{How it is tested.}{A thousand people are drawn at random from the
thousand held-out ones, with replacement, so some appear twice and others not at
all. Both models are scored on exactly that draw and the difference is noted. That
is repeated two thousand times, and the middle 95\,\% of the two thousand
differences is the range printed in the third column. Because both models always
see the same draw, the luck of the sample cancels and only the difference between
the models is left.}

\figpara{How to read a row.}{If the range contains zero, then ``no difference at
all'' is a perfectly ordinary outcome of this experiment, and the two sizes
cannot be told apart. Every row here contains zero, including the last one, which
spans the entire ladder from 5{,}000 to 21{,}000 people and whose best estimate
is 0.04.}

\figpara{What that licenses me to say.}{Not ``the gain from more people is
small''. \emph{No gain was measured at all} above 5{,}000 people. \seleven,
1{,}000 test people, 2{,}000 resamples.}}

\vspace{3pt}
\begin{minipage}{\textwidth}\footnotesize
The fitted scaling law says the same thing: a slope of $+0.18\pt$ per ten-fold
increase in users, and a line that accounts for only 12\,\% of the spread of
the eight points it is fitted to. At that slope one extra point would require
about 1.5 billion users, the absurdity that shows the curve is flat rather
than merely slow. The correct statement is not ``the gain is small'' but
``no gain was measured''.
\end{minipage}
\figsource{\nbfinal}{train\_cellid\_final.ipynb}{the paired-bootstrap cells, 2,000 resamples}
\end{table}


% T10 -- Every metric across the four tiers (slide 78).
\begin{table}[!htbp]
\centering\small
\caption{Six ways of measuring, and all six say the same thing}
\label{tab:all-metrics}
\begin{tabular}{@{}lrrrrrr@{}}
\toprule
\textbf{Way of measuring} & \textbf{1{,}900}$^{\dagger}$ & \textbf{5{,}000} & \textbf{10{,}000} & \textbf{15{,}000} & \textbf{21{,}000} & \textbf{Spread} \\
\midrule
Top-1, the headline         & 52.4\,\% & 54.5\,\% & 54.6\,\% & 54.9\,\% & 54.5\,\% & 0.4\pt \\
Top-3                       & 71.4\,\% & 72.0\,\% & 72.1\,\% & 72.1\,\% & 72.2\,\% & 0.2\pt \\
Top-5                       & 78.6\,\% & 78.9\,\% & 78.9\,\% & 78.9\,\% & 79.3\,\% & 0.4\pt \\
Top-1 on a cold start, 30\,\% of the day & 51.4\,\% & 52.3\,\% & 52.6\,\% & 52.3\,\% & 52.1\,\% & 0.5\pt \\
The middle person           &       & 49.1\,\% & 49.2\,\% & 49.2\,\% & 50.3\,\% & 1.2\pt \\
The worst tenth of people   &       & 14.4\,\% & 13.1\,\% & 14.4\,\% & 13.1\,\% & 1.3\pt \\
\bottomrule
\end{tabular}

\figsource{\nbfinal}{train\_cellid\_final.ipynb}{the six-metric evaluation of the four tiers}
\end{table}


% T8 -- Context fraction at 1,000 and 2,000 users (slide 60).
\begin{table}[!htbp]
\centering\small
\caption{The gain does not come from giving the model a long history}
\label{tab:context-scale}
\begin{tabular}{@{}lrrrrr@{}}
\toprule
\textbf{Past given} & \textbf{Moments} & \textbf{The rule} & \textbf{900 people} & \textbf{1{,}900 people} & \textbf{Difference} \\
\midrule
30\,\% & 5{,}491 & 47.1\,\% & 50.0\,\% \ ($+2.9$) & 51.4\,\% \ ($+4.3$) & $+1.4$ \\
40\,\% & 4{,}703 & 46.8\,\% & 49.9\,\% \ ($+3.1$) & 51.4\,\% \ ($+4.6$) & $+1.5$ \\
50\,\% & 3{,}922 & 46.6\,\% & 50.1\,\% \ ($+3.5$) & 51.5\,\% \ ($+4.9$) & $+1.4$ \\
60\,\% & 3{,}142 & 47.0\,\% & 50.3\,\% \ ($+3.3$) & 51.2\,\% \ ($+4.2$) & $+0.9$ \\
70\,\% & 2{,}359 & 47.3\,\% & 50.5\,\% \ ($+3.2$) & 51.6\,\% \ ($+4.3$) & $+1.1$ \\
80\,\% \textit{(ref.)} & 1{,}570 & 48.2\,\% & 51.0\,\% \ ($+2.8$) & 52.4\,\% \ ($+4.2$) & $+1.4$ \\
90\,\% & 781 & 50.3\,\% & 51.9\,\% \ ($+1.6$) & 53.6\,\% \ ($+3.3$) & $+1.7$ \\
\bottomrule
\end{tabular}

\figtext{%
\figpara{What is being varied.}{Two trained models, one on 900 training people
and one on 1{,}900, both graded on the same 100 held-out people, and each row
hands them a different amount of that person's day as history. ``30\,\%'' means
the model sees the first 30\,\% of the day and answers the rest.}

\figpara{How to read a row.}{\textbf{Moments} is how many graded moments that row
rests on, and it falls as more of the day is given away as history.
\textbf{The rule} is what ``repeat the last antenna'' scores on those moments.
The two model columns give the score, with the distance to the rule in brackets.
\textbf{Difference} is how much the larger training set was worth on that row.}

\figpara{What it shows.}{The 1{,}900-person model is ahead of the rule by 4.3
even when it has seen only 30\,\% of the day, and by 4.2 when it has seen 80\,\%.
So what the model contributes is essentially independent of how much history it
is handed. Whatever it has learned, it is not something that needs a long run-up
to become useful.}

\figpara{And the last column is the point of the table.}{Going from 900 to
1{,}900 training people is worth between 0.9 and 1.7 on every row alike. More
people helps by the same amount whether the model is given a third of somebody's
day or nine tenths of it, which is one more sign that the two levers are
independent. \stwok, with the same 100 held-out people for both sizes.}}

\vspace{3pt}
\begin{minipage}{\textwidth}\footnotesize
Graded over the whole day rather than a cut of it, 7{,}745 moments, the two sizes
score 50.0\,\% and 51.0\,\% against a rule of 46.8\,\%.
\end{minipage}
\figsource{\nbfinal}{train\_cellid\_final.ipynb}{the context-fraction evaluation at the smaller tiers}
\end{table}


\textbf{Table~\ref{tab:bootstrap}} presents bootstrap confidence intervals across cohort scale transitions. For all transitions above 5,000 subscribers in \textbf{Table~\ref{tab:bootstrap}}, 95\% confidence intervals contain zero (e.g., net delta between 5,000 and 21,000 subscribers is $+0.04$ percentage points, 95\% CI $[-0.45, +0.52]$). \textbf{Table~\ref{tab:all-metrics}} confirms saturation across six complementary evaluation metrics (top-1, top-3, top-5, weighted metrics). \textbf{Table~\ref{tab:context-scale}} demonstrates that saturation holds consistently across all prefix context fractions (30\% to 90\%).

\FloatBarrier

\subsection{Sensitivity Analysis: User Typology vs. Data Scale}

Evaluating the 21,000-subscriber model across behavioral cohorts reveals substantial variance: top-1 accuracy reaches 62.8\% for stationary subscribers ($G0$, 22\% of cohort), 53.6\% for regular mobile subscribers ($G1$, 37\%), 43.0\% for evening-mobile subscribers ($G2$, 22\%), and 40.4\% for morning-mobile subscribers ($G3$, 18\%). This establishes a performance spread of \textbf{22.4 percentage points across user types}, compared to a marginal \textbf{0.4 percentage point spread across dataset scale} (from 1,900 to 21,000 subscribers in \textbf{Figure~\ref{fig:scaling-ladder}}).

Subscriber mobility typology exerts a $50\times$ stronger impact on prediction accuracy than scaling training cohort size beyond 2,000 subscribers.

\section{Mechanistic Analysis of Saturation Thresholds}

% F -- The two ends of the scaling curve, and the mechanism proposed for each.
\begin{figure}[!htbp]
\centering
\fitwidth{%
\begin{tikzpicture}[font=\small,
  ttl/.style={anchor=north west,align=left,font=\footnotesize\bfseries,text=clmodel},
  axl/.style={font=\tiny,text=clrule},
  gv/.style={font=\scriptsize\bfseries,text=clgain,anchor=north},
]
% ================= panel (a) : the smallest kind becomes estimable ========
\begin{scope}
\node[ttl,text width=8.0cm] at (0,0.40)
  {(a) Why it climbs: the smallest kind of user crosses a hundred members};
\draw[clrule] (0.60,-0.45) -- (0.60,-3.30) -- (8.30,-3.30);
\node[axl,rotate=90,anchor=south] at (0.36,-1.90)
  {people in the evening-moving user};
\draw[clloss,dashed,thick] (0.60,-2.34) -- (8.20,-2.34);
\foreach \xx/\hh/\lbl in {1.55/0.54/56, 3.45/1.21/126, 5.35/1.88/196, 7.25/2.55/266}{
  \fill[clGtwo] (\xx-0.42,-3.30) rectangle (\xx+0.42,{-3.30+\hh});
  \node[font=\scriptsize,text=clrule,anchor=south] at (\xx,{-3.26+\hh}) {\lbl};
}
\node[anchor=north west,align=left,text width=2.10cm,font=\scriptsize,
      text=clloss] at (0.75,-0.62)
  {about 100 members: enough to estimate a rhythm};
\draw[clloss] (1.70,-1.60) -- (1.70,-2.34);
\foreach \xx/\np/\gg in {1.55/400/{+0.6}, 3.45/900/{+2.8}, 5.35/{1{,}400}/{+3.5},
                         7.25/{1{,}900}/{+4.2}}{
  \node[anchor=north,font=\scriptsize,text=clrule] at (\xx,-3.36) {\np};
  \node[gv] at (\xx,-3.78) {\gg};
}
\node[axl,anchor=north] at (4.45,-4.22)
  {top row: people trained on \ \ \ bottom row: what the model gains over the rule};
\end{scope}

% ================= panel (b) : nothing left to add =======================
\begin{scope}[shift={(8.95,0)}]
\node[ttl,text width=5.9cm] at (0,0.40)
  {(b) Why it stops: there is nothing left to see};
\node[anchor=north west,align=left,text width=6.10cm,font=\scriptsize,
      text=clrule] at (0,-0.42)
  {share of the test people's antennas already seen during training};
% bar 1
\node[anchor=west,font=\scriptsize,text=clmodel] at (0,-1.42) {400 people};
\fill[clnonsig!55] (2.60,-1.54) rectangle (5.90,-1.30);
\fill[clbase]      (2.60,-1.54) rectangle (5.86,-1.30);
\node[anchor=west,font=\scriptsize\bfseries,text=clbase] at (6.00,-1.42) {98.9};
% bar 2
\node[anchor=west,font=\scriptsize,text=clmodel] at (0,-1.97) {5{,}000 and above};
\fill[clbase] (2.60,-2.09) rectangle (5.90,-1.85);
\node[anchor=west,font=\scriptsize\bfseries,text=clbase] at (6.00,-1.97) {100};
\draw[clrule!60] (0,-2.38) -- (6.60,-2.38);
\node[anchor=north west,align=left,text width=6.40cm,font=\scriptsize,
      text=clmodel] at (0,-2.52)
  {And the four shapes of Figure~\ref{fig:group-profiles} are essentially
   identical at every size, so 5{,}000 people already describe this population
   completely.};
\node[anchor=north west,draw=clrule,fill=clsoft,rounded corners=2pt,inner sep=5pt,
      align=left,text width=6.05cm,font=\scriptsize] at (0,-3.62)
  {The sixteen thousand people added after that bring no antenna and no
   behaviour the model has not already seen hundreds of times. They are worth
   \textbf{\textcolor{clgain}{+0.1}}.};
\end{scope}
\end{tikzpicture}}
\caption{Why the curve climbs, and why it stops}
\label{fig:why-it-stops}
\figtext{%
\figpara{Panel (a): one mechanism fits the timing of the jump.}{The bars are how
many people fall in the least populated of the four kinds, the evening-moving
user, at 14\,\% of the population. Under each bar is what the model gained over
the one-line rule at that training size. The count crosses roughly a hundred
members exactly between the two sizes where the gain jumps from 0.6 to 2.8. A
label can only carry information if the group it names has enough inhabitants
for its rhythm to be estimated at all, and below that the model is right to
ignore it, which is precisely what the failed run of Chapter~\ref{ch:groups}
showed it doing. Cross-training did not become a better idea at 900 people than
at 400; the kinds simply became large enough to be worth reading.}

\figpara{Panel (b): above the threshold every new person is a duplicate.}{With
113 masts, four kinds of user and every antenna in the test set already seen
during training at 5{,}000 people, more people can only bring more of the same.
This is an explanation that fits, not a proof: the alternative reading, that the
training budget rather than the population ran out, is stated in
Section~\ref{sec:scale-limits}.}}
\figsource{\nbfinal}{train\_cellid\_final.ipynb}{the eight-tier scaling notebook; the member counts are the group shares of Chapter~\ref{ch:groups} applied to each size}
\end{figure}


\textbf{Figure~\ref{fig:why-it-stops}} illustrates the mechanics underlying performance saturation:
- \textbf{\panelref{fig:why-it-stops}{a}} shows that at 2,000 subscribers, even the smallest behavioral subgroup ($G2$) achieves sufficient trajectory samples ($\approx 280$ subscribers) to robustly estimate subgroup transition matrices.
- \textbf{\panelref{fig:why-it-stops}{b}} shows spatial cell vocabulary coverage. Beyond 2,000 subscribers, 100\% of test cell IDs are present in the training set vocabulary, meaning additional subscribers contribute redundant transition paths across previously observed cell graphs.

\FloatBarrier

\subsection{Methodological Constraints and Scope Boundaries}
\label{sec:scale-limits}

Scaling experiments maintained constant compute step budgets (8,000 gradient update steps per run). At 5,000 subscribers (5,379 training sequences), epoch coverage reaches 1.5 passes. At 21,000 subscribers (22,522 sequences), 8,000 steps correspond to 36\% single-pass epoch coverage. Consequently, reported saturation bounds hold under fixed compute step constraints and fixed parameter adapter capacity.

\section{Evaluation of Trajectory Length Equalization}

An alternative preprocessing approach was evaluated: equalizing subscriber trajectory lengths by oversampling short trajectories to a uniform 200 records per subscriber.

% F15 -- Levelling every user to 200 records: what the score looks like, and what it is.
\begin{figure}[t]
\centering
\replegend{\swatch{clbase}~persistence baseline \qquad \swatch{clmodel}~model}
\fitwidth{%
\begin{tikzpicture}
\begin{groupplot}[
  group style={group size=2 by 1, horizontal sep=2.2cm, ylabels at=edge left},
  repbar, width=0.42\linewidth, height=5.4cm, xtick=data,
  title style={font=\footnotesize,yshift=1pt},
]
% ---- (a) each run against its own baseline ----
\nextgroupplot[title={(a) each run against \emph{its own} baseline},
  symbolic x coords={nat,lev},
  xticklabels={natural lengths,levelled to 200},
  ymin=0, ymax=122, ytick={0,20,...,80},
  ylabel={top-1 accuracy (\%)}, /pgf/bar width=16pt,
  x tick label style={font=\scriptsize}]
\addplot[fill=clbase,draw=none]  coordinates {(nat,43.1)(lev,76.5)};
\addplot[fill=clmodel,draw=none] coordinates {(nat,46.8)(lev,74.9)};
\node[font=\small,text=clgain,anchor=south] at (axis cs:nat,60) {$+3.7\pt$};
\node[font=\small,text=clloss,anchor=south] at (axis cs:lev,102) {$-1.6\pt$};
%
% ---- (b) the arithmetic of copying ----
\nextgroupplot[title={(b) what the 4{,}000 scored positions are},
  symbolic x coords={cop,real},
  xticklabels={{copies\\2{,}262 of 4{,}000},{real decisions\\1{,}738 of 4{,}000}},
  ymin=0, ymax=125, ytick={0,25,...,100},
  /pgf/bar width=16pt,
  x tick label style={font=\scriptsize,align=center}]
\addplot[fill=clbase,draw=none]  coordinates {(cop,100.0)(real,45.9)};
\addplot[fill=clmodel,draw=none] coordinates {(real,42.2)};
\node[font=\scriptsize,text=clloss,anchor=south,xshift=-11pt] at (axis cs:real,58) {$-3.7\pt$};
\end{groupplot}
\end{tikzpicture}}
\caption{Levelling every user to 200 records}
\label{fig:resampling}
\figtext{%
Stretching every sequence to exactly 200 records makes the score look
spectacular and turns the gap negative. (a)~Each run against \emph{its own}
persistence baseline, natural lengths on the left and levelled sequences on
the right; only the comparison inside a pair means anything. The levelled
model reaches 74.9\,\%, twenty-eight points above the natural-length run and
by a wide margin the highest absolute accuracy in the whole project, but the
baseline on those same sequences reaches 76.5\,\% and the gap changes sign,
from $+3.7\pt$ to $-1.6\pt$. (b)~The reason is arithmetic rather than
modelling. Of the 4{,}000 scored positions (100 users $\times$ the last 40 of
their 200 records), 2{,}262 ask for a record copied from the one before, and
any rule that repeats is right on all of them by construction; both the model
and the baseline are handed those free answers, which is why both rose by
about thirty points. On the 1{,}738 positions that are real decisions the
baseline scores 45.9\,\% against the model's 42.2\,\%. The technique does
deliver what it promises, the longest-to-shortest ratio falling from 200:19 to
1:1 and the Gini coefficient of per-user weight to zero, but a copied step
never changes site, so levelling compresses the between-user spread in site-change rate by a factor of $2.3$ and drags every user towards the homebody
profile, which disables the group detection the project's only real gain
depends on. The model bar in panel~(b) is derived from the reported 74.9\,\%,
not re-measured: it assumes the model never misses a copy, and rises to
44.8\,\% if it misses 2\,\% of them, both ends staying below the baseline.}
\end{figure}


\textbf{Figure~\ref{fig:resampling}} evaluates trajectory equalization: \textbf{\panelref{fig:resampling}{a}} shows apparent absolute accuracy gains (reaching 74.9\%), while \textbf{\panelref{fig:resampling}{b}} reveals underlying distortion. As detailed in \textbf{\panelref{fig:resampling}{b}}, oversampling short trajectories fills the evaluation set with 2,262 pure copy moments that inflate the baseline to 76.5\%, while performance on the 1,738 real decision moments falls 3.7 percentage points behind the baseline (42.2\% vs. 45.9\%). Artificially equalizing trajectory lengths distorts empirical transition distributions.

\section{Conclusion}
\label{sec:retro-scale}

Cohort scaling analysis demonstrates that next-cell prediction accuracy saturates at approximately 2,000 subscribers under parameter-efficient fine-tuning (+4.2 percentage points over baseline in \textbf{Figure~\ref{fig:scaling-ladder}}). Expanding training cohorts from 2,000 to 21,000 subscribers yields negligible accuracy gains (+0.04 percentage points, $p > 0.05$ in \textbf{Table~\ref{tab:bootstrap}}). Furthermore, individual subscriber mobility typology accounts for a 22.4 percentage point accuracy spread, demonstrating that individual user characterization exerts a far stronger influence on predictability than raw data scale.

